% ch06 — Geodesic Motion
\chapter{Geodesic Motion}
\label{ch:geodesics}

Geodesics are the trajectories of freely falling particles and photons---the
central objects computed by the ray tracer. This chapter develops both the
Lagrangian and Hamiltonian formulations, derives the constants of motion, and
describes the numerical strategies used in the codebase for integration,
conservation monitoring, and ray termination.

\section{Lagrangian and Hamiltonian Formulations}
\label{sec:geo-formulations}

We derive the geodesic equations from the Euler-Lagrange equations and
recast them in Hamiltonian form using the super-Hamiltonian
$\mathcal{H} = \tfrac{1}{2}\,g^{\mu\nu} p_\mu p_\nu$.

\section{Constants of Motion}
\label{sec:geo-constants}

We identify the conserved energy, angular momentum, and (in Kerr) Carter
constant from Killing vectors and the Killing tensor, reducing the equations
of motion.

\section{Null Geodesics}
\label{sec:geo-null}
\coderef{Geodesic}{Hamiltonian}

We specialise to null geodesics ($\mathcal{H} = 0$), the light rays traced
by the renderer, and discuss the constraint-preserving properties of the
Hamiltonian formulation.

\section{Timelike Geodesics}
\label{sec:geo-timelike}

We treat massive-particle orbits, including circular orbits and their
stability, connecting to accretion disc physics.

\section{Hamiltonian Ray Tracing}
\label{sec:geo-raytracing}

We describe the first-order Hamiltonian system integrated by the codebase,
including the choice of affine parameter and the role of the metric inverse.

\section{Conservation Monitoring}
\label{sec:geo-conservation}

We explain how the codebase monitors the Hamiltonian constraint and conserved
quantities during integration to detect numerical error accumulation.

\section{Termination Conditions}
\label{sec:geo-termination}

We catalogue the stopping criteria---horizon crossing, escape to infinity,
disc intersection---and their implementation.

\section{Photon Rings and Bound Orbits\starmark{}}
\label{sec:geo-photon-rings}

We discuss the structure of bound null orbits around Kerr black holes and
the resulting photon-ring substructure visible in high-resolution images.

\paragraph{Key References.}
The Hamiltonian formulation follows \cite{MTW1973}; the Carter constant and
Kerr geodesics are treated in \cite{Chandrasekhar1983}; photon rings are
discussed in \cite{Johnson2020}.
